% !TEX root = ../main.tex
% !TEX program = lualatex
\documentclass[../main.tex]{subfiles}
\IfSubfilesClassLoaded{\externaldocument{\subfix{../build/main}}}{}
% =============================================================================
% File: 27L_Mentoring_Tips.tex
% =============================================================================
\begin{document}
\IfSubfilesClassLoaded{\chapter{EDO Study Guide}}{}
\section{Mentoring Tips}
\note{Flow cue: this chapter applies the \ac{ldf}; the framework itself is summarized later in Section~\ref{sec:leader-development-framework}. For board answers, describe mentoring as one execution mechanism for EDO character, competence, and culture development.}

\subsection{Summary}
\begin{description}
  \item[\textbf{Definition.}] Mentoring provides personal and professional development advice and
  career counseling; it is not detailing or formal performance counseling~\autocite{edo-5-2-6-mentoring-tips-2025}.
  \item[\textbf{Cadence.}] Formal mentoring is required annually, but informal mentoring is encouraged and
  can occur with any trusted advisor regardless of rank or mentor group~\autocite{edo-5-2-6-mentoring-tips-2025}.
  \item[\textbf{EDO policy.}] The \ac{ldf} sets mentoring expectations and guiding principles focused on
  knowledge transfer, alignment, and long-term community investment~\autocite{edo-5-2-6-mentoring-tips-2025}.
  \item[\textbf{Preparation.}] Effective sessions require read-ahead materials (career planner, \ac{osr}/\ac{psr})
  and clear questions sent in advance to maximize mentor feedback~\autocite{edo-5-2-6-mentoring-tips-2025}.
  \item[\textbf{Follow-through.}] The most important step is follow-up: capture guidance, update plans, and
  maintain contact to build long-term professional networks~\autocite{edo-5-2-6-mentoring-tips-2025}.
\end{description}

\subsection{Mentoring Definition and Steps}
\begin{description}
  \item[\textbf{What mentoring is.}] Mentoring focuses on coaching, development, and career guidance that
  helps \ac{edo}s make informed decisions across professional and personal milestones~\autocite{edo-5-2-6-mentoring-tips-2025}.
  \item[\textbf{What mentoring is not.}] Mentoring does not replace detailing, and it is not solely a review
  of performance counseling or fitness reports~\autocite{edo-5-2-6-mentoring-tips-2025}.
  \item[\textbf{Three-step cycle.}] Prepare, conduct, and follow up; the follow-up step is emphasized as the
  most important for growth and accountability~\autocite{edo-5-2-6-mentoring-tips-2025}.
\end{description}

\subsection{EDO Mentoring Policy and Mentor Groups}
\begin{description}
  \item[\textbf{Leader Development Framework.}] The \ac{ldf} anchors formal mentoring and reinforces
  character, competence, and culture goals for the \ac{edo} community~\autocite{edo-5-2-6-mentoring-tips-2025}.
  \item[\textbf{Guiding principles.}] Mentoring should transfer community knowledge, align professional
  and personal goals, cultivate potential, and invest in future leaders~\autocite{edo-5-2-6-mentoring-tips-2025}.
  \item[\textbf{Mentor groups.}] Mentor groups are sponsored by senior \acp{edo}, led by captains in key
  billets, and administered by a mentor group secretary to track participation and assignments~\autocite{edo-5-2-6-mentoring-tips-2025}.
\end{description}

\subsection{Career Planner Guidance}
\begin{description}
  \item[\textbf{Purpose.}] The career planner helps visualize career paths, communicate preferences, and
  support mentoring discussions with flag officers, mentor group leads, and detailers~\autocite{edo-5-2-6-mentoring-tips-2025}.
  \item[\textbf{Content.}] Maintain several viable paths toward crown jewel jobs or command, and update the
  planner as assignments evolve and mentoring feedback is received~\autocite{edo-5-2-6-mentoring-tips-2025}.
  \item[\textbf{Personal considerations.}] Include geographic preferences and family considerations to ensure
  mentors can provide realistic guidance and options~\autocite{edo-5-2-6-mentoring-tips-2025}.
\end{description}

\subsection{Prepare and Conduct Mentoring}
\begin{description}
  \item[\textbf{Prepare early.}] Provide read-ahead materials at least 48 hours in advance (career planner,
  \ac{osr}/\ac{psr}, and questions) and review your mentor's background to focus the discussion~\autocite{edo-5-2-6-mentoring-tips-2025}.
  \item[\textbf{Core topics.}] Typical discussions include qualification status, assignment planning,
  promotion readiness (\ac{dawia} level, \ac{apm} status), and personal considerations such as family or
  \ac{efm} impacts~\autocite{edo-5-2-6-mentoring-tips-2025}.
  \item[\textbf{Mindset.}] Mentoring advice is autobiographical; be open, take notes, and treat sessions as a
  professional interview and learning opportunity~\autocite{edo-5-2-6-mentoring-tips-2025}.
\end{description}

\subsection{Follow-Up}
\begin{description}
  \item[\textbf{Sustain the relationship.}] Send a thank-you note, share updated planners, and provide
  periodic updates outside formal windows to maintain momentum~\autocite{edo-5-2-6-mentoring-tips-2025}.
  \item[\textbf{Pay it forward.}] Share mentor resources with peers and mentor others as you would like to be
  mentored~\autocite{edo-5-2-6-mentoring-tips-2025}.
\end{description}

\subsection{Quick Review}
\begin{itemize}
  \item Formal mentoring is required annually, but informal mentoring is encouraged throughout the year~\autocite{edo-5-2-6-mentoring-tips-2025}.
  \item The \ac{ldf} and mentor groups guide community standards for mentoring and development~\autocite{edo-5-2-6-mentoring-tips-2025}.
  \item Effective mentoring requires preparation, focused discussion, and disciplined follow-up~\autocite{edo-5-2-6-mentoring-tips-2025}.
  \item Career planners and \ac{osr}/\ac{psr} reviews anchor the session and sharpen promotion readiness~\autocite{edo-5-2-6-mentoring-tips-2025}.
  \item Mentoring works best when advice is internalized and shared forward with others~\autocite{edo-5-2-6-mentoring-tips-2025}.
\end{itemize}

\IfSubfilesClassLoaded{
  \RenewDocumentCommand{\entryname}{}{\textbf{\color{Modern} Acronym}}
  \RenewDocumentCommand{\descriptionname}{}{\textbf{\color{Modern} Definition}}
  \printnoidxglossary[
    type=\acronymtype,
    title=Acronyms,
    style=long-booktabs
  ]}{}
\end{document}
